% Report fragment; requires amsmath, amssymb, booktabs, hyperref.
\section{A second 194-state machine from a symmetric guard}
\label{sec:symmetric194}

The second new machine, $M_{CS}$ or PA194CS, changes the source computation
of Shift and then folds the new reader into an existing decrement control.
It has 194 working states and 388 transition rows, with the same blank
binary one-tape convention as PA194C. Its source has \textbf{3,646 sites}
and \textbf{4,745 primitive outcomes}. The literal table is
\texttt{secondary/pa194CS/machine/pa194CS.tm}; the dense SHA-256 is
\begin{quote}\small
\nolinkurl{1c71002e5b8a2ddd1dd1311a2c54c424764624b611b69c8efbee1ac23c2e6f5f}.
\end{quote}
This is an additional source simplification. It retains the exact cached
command interface of PA194C, while removing an unnecessary restoration
loop from its implementation.

\subsection{The comparison routine is symmetric in its full effect}

The first source instructions of both programs are the following finite
register graph. Here $H=r_5$, $X=r_6$, and a decrement's first successor
is its zero branch.
\[
\begin{array}{r|l}
0&\operatorname{dec}H\;[4,1]\\
1&\operatorname{dec}X\;[3,2]\\
2&\operatorname{jump}0\\
3&\operatorname{jump}5\\
4&\operatorname{dec}X\;[10,5]\\
5&\operatorname{dec}H\;[7,6]\\
6&\operatorname{jump}5\\
7&\operatorname{dec}X\;[9,8]\\
8&\operatorname{jump}5\\
9&\operatorname{jump}11\\
10&\mathsf{halt}.
\end{array}
\]
From any natural pair $(H,X)=(a,b)$ this routine halts exactly when $a=b$.
If $a\ne b$, it reaches site 11 with $H=X=0$, preserving every other
register. To prove this, the loop $0,1,2$ removes one from each positive
component. If $H$ first becomes zero, instruction 4 distinguishes equality
from a positive remainder in $X$. If $X$ first becomes zero while $H$ is
positive, instruction 3 enters cleanup. Sites $5,6$ clear the remaining
$H$, and sites $7,8$ clear the remaining $X$; the latter loop may revisit
the already-zero $H$ and still terminates. Induction on $a+b$ establishes
these statements in PA.

Thus exchanging $(a,b)$ with $(b,a)$ preserves both halting and the entire
surviving register store. Equality of the Boolean test alone would not
suffice for this optimization; clearing \emph{both} comparison registers
is the property needed by the following stack decoder.

\subsection{The shortened Shift}

At the inherited Shift boundary let $(T,H,X)=(t,h,0)$. PA194C performs
\[
\begin{aligned}
(t,h,0)&\longmapsto(0,h,0)\\
       &\longmapsto(h,0,h)\\
       &\longmapsto(h,h,0)\\
       &\longmapsto(h,h,1),
\end{aligned}
\]
then enters the comparison. PA194CS instead performs only
\[
       (t,h,0)\longmapsto(0,h,0)
          \longmapsto(h,0,h)\longmapsto(h,1,h)
\]
and enters that same comparison. The deleted loop restores $X$ into $H$;
the final increment has changed from $X$ to $H$. By the preceding lemma,
both computations halt exactly when $h=1$. Otherwise both reach the
unchanged stack decoder at site 11 with
\[
                         (T,H,X)=(h,0,0).
\]
The unchanged decoder therefore produces the same command store in both
machines. This holds for every natural $t,h$, including $h=1$; there is no
consistency assumption in the argument.

In fact, the source-block equivalence is stronger than the command-boundary
claim. For arbitrary incoming scratch value $x$, the common clear/copy
prefix sends $(T,H,X)=(t,h,x)$ to $(h,0,k)$ with $k=h+x$. The old tail
and guard then compare $(H,X)=(k,1)$, while the new ones compare
$(1,k)$. They again have identical halting and complete surviving-register
effects. Thus this local source proof requires no zero-scratch hypothesis.

The finite source binding is exact. The first 3,644 resolved instructions
of the new source agree with those of the pinned cached controller.
The old final five instructions are
\[
\begin{array}{r|l}
3644&\operatorname{dec}X\;[3647,3645]\\
3645&\operatorname{inc}H\to3646\\
3646&\operatorname{jump}3644\\
3647&\operatorname{inc}X\to3648\\
3648&\operatorname{jump}0,
\end{array}
\]
whereas the new tail is
\[
             3644:\operatorname{inc}H\to3645,
             \qquad3645:\operatorname{jump}0.
\]
The preceding transfer loop's zero exit remains site 3644 and now enters
the shorter tail. An incoming-edge check shows that no instruction outside
the replaced tail enters its interior: the only external edge to the tail
is the zero branch from site 3640, which establishes $H=0$. Consequently
the common source can be partitioned into unchanged single steps and
complete shortened-Shift/comparator blocks. The blocks terminate on every
natural register store, have the same halt outcome or the same full store
at site 11, and have positive duration. Induction on these blocks gives
halting equivalence of the two source programs from site 25 for arbitrary
initial registers. This is a direct local source simulation; the entire
source graphs are not asserted equal.

All AX formulas, command selectors, MP code and scheduler instructions
are retained. The logical registers agree at command
boundaries; physical spare lists need not have equal lengths because the
number of increments changes. Each machine's complete-frame simulation
tracks its own spare-list effects, so this difference causes no omission
in the tape argument.

\subsection{Recovering the state cost by an exact contextual fold}

A direct physical implementation uses one new reader, native control 330:
\[
 (330,0):0R63,\qquad(330,1):1L73.
\]
Relative to PA194C, redirect $(214,1)$ to $1R30$ and $(310,1)$ to
$1R330$. This produces a 195-state intermediate for the shortened source.
Control 41, with actions $0R24$ and $0R25$, is entered only through
$(28,0):0R41$ or $(70,0):0R41$. The new 330 is entered only through
$(310,1):1R330$. Consequently the preceding tape bit distinguishes the
two controls: zero for 41 and one for 330.

Identify $41$ with $330$, replace the retained rows by
\[
                  (330,0):0L31,\qquad(330,1):1L46,
\]
and redirect all old targets 41 to 330. The following exact three-step
paths reproduce the old one-step actions. The intermediate head visits
only the current and preceding tape cells.
\[
\begin{array}{cc|c|l}
\text{preceding}&\text{current}&\text{old action}&\text{controls along replacement}\\\hline
0&0&0R24&330\to31\to21\to24\\
0&1&0R25&330\to46\to61\to25\\
1&0&0R63&330\to31\to47\to63\\
1&1&1L73&330\to46\to62\to73
\end{array}
\]
For every row, the preceding bit and all exterior tape are preserved,
and the final current bit, head coordinate and target agree with the
old action. These four cases and every incoming row are checked against
the \emph{final} quotient table. Ordinary transitions commute with the
projection. Induction gives an all-tape positive simulation from the
specified start; the usual determinism argument gives the converse
halting implication. The original seven AX folds remain separately
verified. Counting from the fully unfolded control table gives
\[
                         135+67-7-1=194.
\]

\subsection{Consistency calibration and verification scope}

The changed Shift has exactly the cached command semantics already
proved for PA194C. Its MP and AX commands are unchanged. Consequently
the same ten-entry proof compiler, all-store universal refutation word,
proof-store soundness invariant and exhaustive scheduler apply without
changing a stack exponent or an axiom schema. Induction on command
executions gives the same source refutation criterion. The guard proof
above uses only natural-number loop induction; the finite source binding
and contextual quotient use primitive recursive trace transformations.
All are formalizable in PA.

The completed literal table is checked on all 4,745 universal source
outcomes, with arbitrary natural primary-register values, both decrement
branches, an arbitrary gap and an arbitrary finite nonempty spare list.
The independent checker rederives the same four list-sweep lemmas from
literal transitions, resolves the new source graph twice, and verifies
the exact shorter tail against the pinned cached source. It does not
assert source-graph identity with PA195C, which would be false.

The \emph{folded} machine reaches the exact complete clean source frame
after \textbf{57,768} blank-tape transitions. The 195-state intermediate
takes 57,594; those counts refer to different tables. Both reach site 25,
native control 84, head $-2$, logical seed $(2,0,0,0,0,0,0,0,0)$ and the
recorded complete spare list. Thus startup, source semantics and full
physical simulation all refer to the delivered 194-state table.

With the same explicitly inherited ordinary syntactic calibration of the
separate-head numerical calculus, these facts give
\[
 \mathrm{PA}\vdash
 \bigl(\neg H_{M_{CS}}\leftrightarrow\operatorname{Con}(\mathrm{PA})\bigr).
\]
The new contribution is the symmetric-guard source proof and its exact
194-state implementation. The logical PA calibration is inherited;
no new end-to-end proof-assistant kernel theorem is asserted.
